% !TeX root = erm_textbook_style_guide.tex
% ERM Textbook Style Guide — formatted with erm_textbook_style.sty
\documentclass[11pt,openany]{book}
\usepackage[utf8]{inputenc}
\usepackage{amsmath,amssymb}
\makeatletter
\def\input@path{{./}{second_ed/style/}}
\makeatother
\usepackage{erm_textbook_style}
\usepackage{listings}
\usepackage{hanging}
\providecommand{\tightlist}{\setlength{\itemsep}{0pt}\setlength{\parskip}{0pt}}
\emergencystretch=2em

% ---------- Code listing style (Iowa palette) ----------
\lstdefinestyle{erm}{
  language=[LaTeX]TeX,
  basicstyle=\ttfamily\small,
  keywordstyle=\color{ERMDarkGold}\bfseries,
  commentstyle=\color{ERMMidGray}\itshape,
  stringstyle=\color{ERMDarkGray},
  backgroundcolor=\color{ERMBoxGray},
  frame=single,
  framesep=4pt,
  framerule=0.4pt,
  rulecolor=\color{ERMLightGray},
  xleftmargin=8pt,
  xrightmargin=4pt,
  breaklines=true,
  breakatwhitespace=true,
  showstringspaces=false,
  tabsize=2,
  aboveskip=0.6\baselineskip,
  belowskip=0.6\baselineskip
}
\lstset{style=erm}

% Inline code: \code{text}
\newcommand{\code}[1]{\texttt{\small #1}}

\begin{document}

\setcounter{chapter}{-1}   % makes this Chapter 0 — a reference/style document
\chapter{ERM Textbook Style Guide}
\label{chap:style-guide}

\chapteroverview{%
This document is the authoritative reference for formatting ERM textbook chapters. It covers the \LaTeX\ package, all custom environments and commands, color palette, typography conventions, and the pre-submission checklist. It is also a live demonstration: every environment described here is used at least once in this document, so the rendered PDF serves as a visual specimen sheet.

Keep this document open alongside any chapter draft. When in doubt about how something should look, find the relevant section here and copy the example code.%
}

% -----------------------------------------------------------------------
\section{Document Setup}
\label{sec:setup}
% -----------------------------------------------------------------------

Every chapter file begins with the same preamble. Copy it exactly; do not improvise.

\begin{lstlisting}
% !TeX root = Chapter_XX_Title.tex
\documentclass[11pt,openany]{book}
\usepackage{graphicx}
\usepackage{amsmath,amssymb}
\makeatletter
\def\input@path{{second_ed/style/}{../style/}}
\makeatother
\usepackage{erm_textbook_style}
\providecommand{\tightlist}{\setlength{\itemsep}{0pt}\setlength{\parskip}{0pt}}
\emergencystretch=2em
\end{lstlisting}

The \code{\textbackslash input@path} declaration lets the style package be found whether you compile from the repository root or from inside the chapter folder. Do not change the document class: \code{book} with \code{11pt} and \code{openany} is required so running heads, chapter numbering, and page layout work correctly. If the chapter has figures generated by R scripts, add a \code{\textbackslash graphicspath} line in the preamble:

\begin{lstlisting}
\graphicspath{{second_ed/rscripts/chapterN/}{../rscripts/chapterN/}}
\end{lstlisting}

When compiling a standalone chapter, set the chapter counter before \code{\textbackslash chapter\{\}} so the number is correct:

\begin{lstlisting}
\setcounter{chapter}{1}   % produces "Chapter 2" when \chapter{} increments it
\chapter{Title}
\end{lstlisting}


% -----------------------------------------------------------------------
\section{Color Palette}
\label{sec:colors}
% -----------------------------------------------------------------------

All colors are defined in \code{erm\_textbook\_style.sty}. Use the named constants in any chapter file; never hard-code hex values.

\begin{ermtable}{Color palette — named constants and their intended uses}{tab:colors}
\begin{tabularx}{\linewidth}{llY}
\toprule
\textbf{Name} & \textbf{Hex} & \textbf{Use} \\
\midrule
\code{ERMIowaBlack}    & \#000000 & Body text, chapter titles, headings \\
\code{ERMIowaGold}     & \#FFCD00 & Chapter title rule, accent highlights \\
\code{ERMDarkGold}     & \#B38600 & Section headings, links, list bullets \\
\code{ERMDarkGray}     & \#3C3C3C & Subsection headings, running head text \\
\code{ERMMidGray}      & \#777777 & Captions, figure notes, source lines \\
\code{ERMLightGray}    & \#E6E6E6 & Running head rule \\
\code{ERMBackground}   & \#FBFAF5 & Chapter overview and managerial insight background \\
\code{ERMBoxGold}      & \#FFF3BF & Learning objectives, worked examples, boardroom questions \\
\code{ERMBoxGray}      & \#F2F2F2 & Regulatory notes, framework summaries \\
\code{ERMBoxBlue}      & \#EAF2F8 & Corporate examples (case studies) \\
\code{ERMBoxGreen}     & \#EAF6EA & Checklists \\
\code{ERMBoxRed}       & \#F9EAEA & Common pitfalls \\
\bottomrule
\end{tabularx}
\tablenote{These colors match the R theme in \code{rscripts/erm\_textbook\_style.r}. Figures produced by that script will share the same palette as the surrounding typeset text.}
\end{ermtable}


% -----------------------------------------------------------------------
\section{Chapter-Opening Structure}
\label{sec:opening}
% -----------------------------------------------------------------------

Every chapter follows this opening sequence, in this order: chapter title → learning objectives → chapter overview → opening corporate example (if applicable) → first section. Do not reorder these elements.

\subsection{Chapter Title and Label}

\begin{lstlisting}
\setcounter{chapter}{N-1}
\chapter{Title of the Chapter}
\label{chap:short-label}
\end{lstlisting}

The \code{\textbackslash chapter} command produces the ``Chapter~N'' prefix and the Iowa Gold rule automatically. Do not add numbering or decorative rules manually.

\subsection{Learning Objectives}

\begin{lstlisting}
\begin{learningobjectives}
  \item Verb-first objective stating a specific skill.
  \item Second objective.
\end{learningobjectives}
\end{lstlisting}

Yellow-gold background (\code{ERMBoxGold}), dark gold border. Each item should begin with an action verb: Explain, Compare, Identify, Apply, Describe, Evaluate, Calculate, Construct. Aim for 5--7 items. Fewer than 4 is thin; more than 8 suggests the chapter is trying to do too much.

\subsection{Chapter Overview}

\begin{lstlisting}
\chapteroverview{%
  First paragraph of the overview.

  Second paragraph if needed.%
}
\end{lstlisting}

Black border with a white title bar on a cream background. Write 2--3 paragraphs: what the chapter covers, why it matters, and how it connects to adjacent chapters. Write this section last, after the chapter content is substantially complete.

\subsection{Opening Corporate Example}

\begin{lstlisting}
\begin{corporateexample}[Company Name and the Event Title]
  Narrative prose...
  \source{Citation for the primary source.}
\end{corporateexample}
\end{lstlisting}

Blue background (\code{ERMBoxBlue}), dark gray border, auto-numbered ``Corporate Example~N.M: Title.'' Place immediately after \code{\textbackslash chapteroverview}, before the first \code{\textbackslash section}. End with \code{\textbackslash source\{\}} citing the primary source. Write as narrative prose; do not use bullet lists inside corporate example boxes.


% -----------------------------------------------------------------------
\section{Section Hierarchy}
\label{sec:sections}
% -----------------------------------------------------------------------

\begin{lstlisting}
\section{Section Title}        % large bold black, numbered
\subsection{Subsection Title}  % bold dark gray, numbered
\subsubsection{Sub-sub title}  % bold dark gray, numbered -- use sparingly
\end{lstlisting}

Sections are numbered automatically. Do not add numbers manually. Use \code{\textbackslash section*\{\}} only for unnumbered end-matter (References, Key Terms, etc.) and always pair it with \code{\textbackslash addcontentsline\{toc\}\{section\}\{Title\}} so the entry appears in the table of contents.

\textbf{Inline sub-headings} within a subsection use \code{\textbackslash textbf\{Heading\}} on its own paragraph line, followed immediately by the prose. This is preferable to a \code{\textbackslash subsubsection} when the heading is a label or category rather than a structural section break.

\begin{lstlisting}
\textbf{The Friction: Something Specific}

Prose follows here. At least three sentences are required
(see Section~\ref{sec:paragraphs}).
\end{lstlisting}


% -----------------------------------------------------------------------
\section{Box Environments}
\label{sec:boxes}
% -----------------------------------------------------------------------

All boxes use the shared \code{erm box} tcolorbox style: auto-numbered within the chapter, breakable across pages, consistent padding and corner radius. Each type has a distinct color scheme that signals its purpose to readers.

\subsection{Summary of Available Box Types}

\begin{ermtable}{Box environments — color schemes and intended uses}{tab:boxes}
\begin{tabularx}{\linewidth}{lllY}
\toprule
\textbf{Environment} & \textbf{Background} & \textbf{Border} & \textbf{Use for} \\
\midrule
\code{corporateexample}  & Blue       & Dark gray   & Real corporate case studies \\
\code{workedexample}     & Gold       & Dark gold   & Solved quantitative problems \\
\code{regulatorynote}    & Gray       & Mid gray    & Laws, regulations, frameworks, standards \\
\code{managerialinsight} & Cream      & Black left rule & Practitioner perspective, ``in practice'' notes \\
\code{commonpitfall}     & Red        & Dark gray   & Mistakes to avoid \\
\code{whythismatters}    & Cream      & Gold left rule & Motivating framing before a technical section \\
\code{boardroomquestion} & Gold       & Black       & Questions a board or CRO should ask \\
\code{checklist}         & Green      & Dark gray   & Step-by-step implementation checklists \\
\bottomrule
\end{tabularx}
\end{ermtable}

\subsection{Syntax}

All numbered box environments take an optional title argument in square brackets:

\begin{lstlisting}
\begin{corporateexample}[Optional Title After the Colon]
  Content...
\end{corporateexample}

\begin{regulatorynote}[Optional Title]
  Content...
\end{regulatorynote}
\end{lstlisting}

The \code{checklist} environment takes an optional title that replaces the default ``Checklist'':

\begin{lstlisting}
\begin{checklist}[ERM Implementation Checklist]
  \item First item.
  \item Second item.
\end{checklist}
\end{lstlisting}

\subsection{Source and Note Lines Inside Boxes}

Two commands produce small italic gray lines at the foot of a box or figure:

\begin{lstlisting}
\source{Full citation for the source of the material.}
\figurenote{Methodological note or explanation.}
\end{lstlisting}

\code{\textbackslash source} produces \textit{Source.}~text; \code{\textbackslash figurenote} produces \textit{Note.}~text. Use both inside boxes and beneath figures as appropriate.

\subsection{Box Placement Guidelines}

Place boxes in the flow of the text where they support the surrounding prose, not as afterthoughts at the end of a section. Each box must earn its place: if the content reads naturally as prose, write it as prose. The \code{corporateexample} environment typically appears once per chapter, at the opening, but a second case study mid-chapter is acceptable if it genuinely illustrates a distinct point. The \code{workedexample} environment should follow the theory it illustrates immediately, not be deferred to the end of the section. Do not use more than one \code{commonpitfall} or \code{boardroomquestion} per chapter unless the chapter specifically focuses on implementation.

\subsection{Live Specimens}

The following boxes demonstrate each environment at a glance.

\begin{workedexample}[Computing Expected Tax Savings from Hedging]
A firm faces a progressive tax schedule: 20\% on income below \$100M, 30\% above. Without hedging, income alternates between \$50M and \$150M. Expected annual tax is \$22.5M. With a hedge locking income at \$100M, expected tax falls to \$20M. Annual tax saving: \$2.5M.
\end{workedexample}

\begin{managerialinsight}[What Risk Ownership Means in Practice]
Assigning a risk owner is not the same as assigning blame. The risk owner is accountable for monitoring the risk, keeping the risk register entry current, and escalating when circumstances change materially. They are not expected to control every factor that affects the risk --- only to ensure that someone is always watching it and that the organization does not drift, as GM did, because no one was assigned the enterprise-level question.
\end{managerialinsight}

\begin{commonpitfall}[Treating the Risk Register as a Compliance Document]
A risk register that is updated once a year, at budget time, to satisfy an audit requirement is not risk management. It is documentation theater. The register should be a live working document reviewed at least quarterly, with risk owners actively updating severity assessments as conditions change. If nobody has touched a risk entry in six months, it almost certainly needs attention.
\end{commonpitfall}

\begin{whythismatters}[Why Risk Appetite Statements Fail]
Most organizations write risk appetite statements. Fewer than half of those statements actually influence decisions. The common failure is writing appetite at a level of generality so high --- ``we seek moderate risk'' --- that no one can use it to evaluate a specific proposal. A useful appetite statement specifies what the organization will not do, names the metrics that bound acceptable exposure, and is reviewed by the board when strategy changes. Vague statements protect no one.
\end{whythismatters}

\begin{boardroomquestion}[Governance of Emerging Risks]
When was the last time the board received a report on risks that did not yet appear on the enterprise risk register? What is the process for identifying risks that are material but not yet visible in current operations?
\end{boardroomquestion}

\begin{checklist}[Pre-Submission Chapter Checklist]
  \item Document class is \code{\textbackslash documentclass[11pt,openany]\{book\}}.
  \item \code{\textbackslash usepackage\{erm\_textbook\_style\}} is present and \code{\textbackslash input@path} is set.
  \item \code{\textbackslash chapter\{Title\}} with \code{\textbackslash label\{chap:...\}} is present.
  \item \code{learningobjectives} environment is present (5--7 items).
  \item \code{\textbackslash chapteroverview\{...\}} is present.
  \item Opening \code{corporateexample} box has a \code{\textbackslash source\{...\}} line.
  \item All boxes use the correct named environment (see Table~\ref{tab:boxes}).
  \item \code{keytakeaways} environment is present; items mirror the learning objectives.
  \item All figures use \code{theme\_erm\_textbook()} and are saved at $\geq$300~dpi.
  \item \code{\textbackslash section*\{References\}} with \code{\textbackslash addcontentsline} is present.
  \item Compiled twice with \code{pdflatex}; no errors in the log.
  \item Auxiliary files (\code{.aux}, \code{.log}, etc.) deleted before committing.
\end{checklist}


% -----------------------------------------------------------------------
\section{End-of-Chapter Elements}
\label{sec:endmatter}
% -----------------------------------------------------------------------

End-of-chapter material appears in this fixed order: Key Takeaways → Key Terms (if present) → Short-Answer Questions (if present) → Discussion Questions (if present) → Mini-Case (if present) → Further Reading (if present) → References. Do not reorder these elements.

\subsection{Key Takeaways}

\begin{lstlisting}
\begin{keytakeaways}
\item
  \textbf{Bold Label}: One or more sentences describing the takeaway.
\item
  \textbf{Second Label}: Sentence.
\end{keytakeaways}
\end{lstlisting}

Produces an unnumbered \code{\textbackslash section*\{Key Takeaways\}} heading with a TOC entry automatically. Each item has a bold label followed by a colon and the content. Aim for 5--8 items. Each learning objective stated at the opening should have a corresponding takeaway here.

\subsection{Key Terms}

\begin{lstlisting}
\begin{keyterms}
  \item[\keyterm{Enterprise risk management}] Definition in full sentences.
  \item[\keyterm{Risk appetite}] Definition.
\end{keyterms}
\end{lstlisting}

Use \code{\textbackslash keyterm\{term\}} to bold the term. The environment produces a hanging-indent description list.

\subsection{Discussion Questions and Short-Answer Questions}

\begin{lstlisting}
\begin{discussionquestions}
\item Open-ended question requiring synthesis?
\end{discussionquestions}

\begin{shortanswerquestions}
\item Question with a specific correct answer.
\end{shortanswerquestions}
\end{lstlisting}

\subsection{Mini-Case}

\begin{lstlisting}
\begin{minicase}{Company Name or Scenario Title}
  Narrative text for the case. Two to four paragraphs minimum.

  \begin{discussionquestions}
  \item Case-specific question.
  \end{discussionquestions}
\end{minicase}
\end{lstlisting}

\subsection{Further Reading}

\begin{lstlisting}
\begin{furtherreading}
\item Author (Year). \emph{Title}. Publisher.
\end{furtherreading}
\end{lstlisting}

\subsection{References}

\begin{lstlisting}
\section*{References}
\addcontentsline{toc}{section}{References}
\begingroup
\sloppy

Author, A. (Year). \emph{Title}. Publisher.

\endgroup
\end{lstlisting}

Use \code{\textbackslash begingroup \textbackslash sloppy \ldots \textbackslash endgroup} to prevent overfull hboxes in long URLs. Format follows APA~7th edition. Journal titles and book titles use \code{\textbackslash emph\{\}}. Wrap URLs in \code{\textbackslash url\{\}}.


% -----------------------------------------------------------------------
\section{Figures and Tables}
\label{sec:figures-tables}
% -----------------------------------------------------------------------

\subsection{Figures}

\begin{lstlisting}
\begin{figure}[htbp]
  \centering
  \includegraphics[width=\textwidth]{filename.png}
  \caption{Caption sentence describing what the figure shows.}
  \label{fig:short-label}
  \figurenote{Methodological note, e.g., indexing base or data source.}
\end{figure}
\end{lstlisting}

All figures produced by R must use \code{theme\_erm\_textbook()} from \code{rscripts/erm\_textbook\_style.r}. Save as PNG at 300~dpi minimum, or as PDF for vector output. Every figure must be referenced in the running text before it appears: ``Figure~\code{\textbackslash ref\{fig:label\}} shows\ldots''.

\subsection{Tables}

\begin{lstlisting}
\begin{ermtable}{Caption text here}{tab:label}
  \begin{tabularx}{\linewidth}{lYY}
    \toprule
    Column 1 & Column 2 & Column 3 \\
    \midrule
    Row data & Data     & Data     \\
    \bottomrule
  \end{tabularx}
  \tablenote{Table source or methodological note.}
\end{ermtable}
\end{lstlisting}

The \code{Y} column type (defined as \code{\textbackslash raggedright\textbackslash arraybackslash X}) is used for text-heavy columns. Use standard \code{l}, \code{c}, \code{r} for short or numeric columns. End tables with \code{\textbackslash tablenote\{\}} for source attribution.


% -----------------------------------------------------------------------
\section{Mathematics}
\label{sec:math}
% -----------------------------------------------------------------------

Inline math uses \code{\textbackslash( \ldots \textbackslash)}. Display math uses \code{\textbackslash[ \ldots \textbackslash]} for unnumbered equations and the \code{equation} environment for numbered ones.

\begin{lstlisting}
% Inline
The expected return is \( E[R_i] = R_f + \beta_i (E[R_M] - R_f) \).

% Display, unnumbered
\[
\text{Cost of Risk} = \text{Expected Retained Losses}
                    + \text{Risk Control Costs} + \ldots
\]

% Display, numbered
\begin{equation}
  \label{eq:cost-of-risk}
  C = L + K + F + I
\end{equation}
\end{lstlisting}

Reference numbered equations as Equation~\eqref{eq:label}. Do not use \code{\$\ldots\$} for inline math or \code{\$\$\ldots\$\$} for display math; both are deprecated in \LaTeXe.


% -----------------------------------------------------------------------
\section{Paragraph Length}
\label{sec:paragraphs}
% -----------------------------------------------------------------------

Paragraphs in body text must contain at least three sentences. One- and two-sentence paragraphs are prohibited.

This is a textbook, not a trade book or a slide deck. Students read chapters sequentially and need developed argument, not fragments. A paragraph that ends after two sentences almost always means one of two things: the idea was not fully explained, or it belongs merged with the paragraph before or after it. Either way, the fix is to write more, not to let it stand as-is.

The only exceptions to this rule are: the final sentence of a section used deliberately as a one-line conclusion (use sparingly, no more than once per section); items inside \code{keytakeaways}, \code{learningobjectives}, or other list environments, where each item is structurally a single point; and captions, source lines, and figure notes. When reviewing a draft, flag every paragraph of fewer than three sentences and either expand it or absorb it into an adjacent paragraph.


% -----------------------------------------------------------------------
\section{Emphasis, Bold, and Key Terms}
\label{sec:emphasis}
% -----------------------------------------------------------------------

Use the following conventions consistently:

\begin{itemize}
  \item \code{\textbackslash emph\{text\}} --- italics for emphasis, for titles of works mentioned inline, and for technical terms on their first use in a chapter.
  \item \code{\textbackslash textbf\{text\}} --- bold for inline sub-headings and for bold labels in \code{keytakeaways} items.
  \item \code{\textbackslash keyterm\{text\}} --- bold; use exclusively inside the \code{keyterms} environment for formally defined terms.
  \item Do not underline. Do not use \textsc{small caps} or ALL CAPS for emphasis.
\end{itemize}


% -----------------------------------------------------------------------
\section{Lists}
\label{sec:lists}
% -----------------------------------------------------------------------

Lists use Iowa Dark Gold bullets automatically via the style package. No additional configuration is needed:

\begin{lstlisting}
\begin{itemize}
  \item First item.
  \item Second item.
\end{itemize}

\begin{enumerate}
  \item First numbered step.
  \item Second numbered step.
\end{enumerate}
\end{lstlisting}

Use \code{\textbackslash tightlist} inside boxes or anywhere vertical spacing should be compressed. Use a list only when items are genuinely discrete and parallel. Three or fewer items almost always read better as prose (``X, Y, and~Z''). Do not convert every set of related points into a bullet list.


% -----------------------------------------------------------------------
\section{Cross-References}
\label{sec:crossrefs}
% -----------------------------------------------------------------------

Always use \code{\textbackslash label} and \code{\textbackslash ref} / \code{\textbackslash eqref} rather than hard-coded numbers.

\begin{lstlisting}
as shown in Figure~\ref{fig:jnj-tylenol}
see Section~\ref{sec:erm-evolution}
Chapter~\ref{chap:why-risk-management-matters}
Equation~\eqref{eq:cost-of-risk}
\end{lstlisting}

Use \code{\textasciitilde} (non-breaking space) between the word and the \code{\textbackslash ref} to prevent line breaks. Follow this labeling convention:

\begin{ermtable}{Label prefixes by element type}{tab:labels}
\begin{tabularx}{\linewidth}{llY}
\toprule
\textbf{Type} & \textbf{Prefix} & \textbf{Example} \\
\midrule
Chapter    & \code{chap:} & \code{\textbackslash label\{chap:enterprise-risk-management\}} \\
Section    & \code{sec:}  & \code{\textbackslash label\{sec:erm-evolution\}} \\
Subsection & \code{subsec:} & \code{\textbackslash label\{subsec:coso\}} \\
Figure     & \code{fig:}  & \code{\textbackslash label\{fig:jnj-tylenol\}} \\
Table      & \code{tab:}  & \code{\textbackslash label\{tab:coso-components\}} \\
Equation   & \code{eq:}   & \code{\textbackslash label\{eq:cost-of-risk\}} \\
\bottomrule
\end{tabularx}
\end{ermtable}


% -----------------------------------------------------------------------
\section{Footnotes}
\label{sec:footnotes}
% -----------------------------------------------------------------------

Use standard \LaTeX\ footnotes for brief clarifying information that would interrupt the main text:

\begin{lstlisting}
Text in the chapter.\footnote{Explanatory note that does not fit inline.}
\end{lstlisting}

Keep footnotes to one or two sentences. If a footnote is growing toward a paragraph, it either belongs in the running text or should be cut. Do not use footnotes for citations; those belong in the References section.


% -----------------------------------------------------------------------
\section{File Organization}
\label{sec:files}
% -----------------------------------------------------------------------

Each chapter lives in its own numbered folder. R figures go in \code{rscripts/chapterN/} subfolders and are referenced via \code{\textbackslash graphicspath} in the chapter preamble.

\begin{lstlisting}[language={}]
second_ed/
  chapter 1/
    Chapter_01_Why_Risk_Management_Matters.tex
    Chapter_01_Why_Risk_Management_Matters.md
    intro c1.md
  chapter 2/
    Chapter_02_Enterprise_Risk_Management.tex
    Chapter_02_Enterprise_Risk_Management.md
    intro c2.md
  chapter 3/
    Chapter_03_Risk_Identification.tex
    Chapter_03_Risk_Identification.md
    intro c3.md
  style/
    erm_textbook_style.sty
    erm_textbook_style_guide.tex   % <-- this document
    erm_textbook_style_guide.md
    erm_textbook_style_guide.pdf
  rscripts/
    erm_textbook_style.r
    chapter1/
      jnj_tylenol/
        c1_jnjplot.png
    chapterN/
\end{lstlisting}

Keep all auxiliary \LaTeX\ files (\code{.aux}, \code{.log}, \code{.out}, \code{.toc}, etc.) out of version control. Add them to \code{.gitignore}.


% -----------------------------------------------------------------------
\section{Compilation}
\label{sec:compilation}
% -----------------------------------------------------------------------

Compile with \code{pdflatex} from the chapter folder or the repository root. Two passes are required to resolve cross-references and TOC entries:

\begin{lstlisting}[language={}]
pdflatex -interaction=nonstopmode Chapter_XX_Title.tex
pdflatex -interaction=nonstopmode Chapter_XX_Title.tex
\end{lstlisting}

If \code{biber} is ever added for bibliography management, the sequence becomes \code{pdflatex} $\to$ \code{biber} $\to$ \code{pdflatex} $\to$ \code{pdflatex}. Currently, references are typed manually, so \code{biber} is not needed.

% -----------------------------------------------------------------------
\section{Chapter-Specific Notes}
\label{sec:chapter-notes}
% -----------------------------------------------------------------------

These notes record structural and stylistic decisions made in completed chapters. They govern how later chapters should handle the same concepts.

\subsection*{Chapter 1 --- Why Risk Management Matters}

Opening case: Johnson \& Johnson Tylenol recall (\code{corporateexample}). Establishes the core Modigliani--Miller irrelevance-propositions framework and the value creation argument for hedging. The Tylenol case motivates hedging under financial distress; the chapter is primarily theoretical with one applied box.

\subsection*{Chapter 2 --- Enterprise Risk Management}

Opening case: GM ignition switch crisis (\code{corporateexample}). The GM case is the motivating example for the entire ERM integration argument---siloed risk management versus enterprise-level oversight. The COSO framework is summarized in a \code{regulatorynote} box. Convention established here: \code{---} separators between major sections; opening narrative flows as prose before the first numbered section.

\subsection*{Chapter 3 --- Risk Identification}

Opening case: Netflix and the Red Envelope Problem (\code{corporateexample}). Two conventions introduced here that carry forward to all subsequent chapters:

\textbf{Enterprise-before-silo framing.} The chapter overview must state explicitly that the goal of risk identification is not four separate lists but a single integrated view. Do not let taxonomy sections drift into functional silos.

\textbf{Definition + vivid example at quadrant openings.} Each risk quadrant section opens with a short concrete example immediately after the definition paragraph, before moving to characteristics and management approaches. Canonical examples:
\begin{itemize}[leftmargin=1.5em]
  \item \emph{Hazard}: a warehouse fire that stops shipments overnight
  \item \emph{Financial}: a sudden euro appreciation that makes every imported component more expensive the moment the exchange rate moves
  \item \emph{Operational}: a payment processing crash during peak holiday traffic that prevents purchases for six hours
  \item \emph{Strategic}: a dense physical store network that was a competitive advantage in 1995 and a strategic risk a decade later
\end{itemize}

\textbf{Business interruption classification.} Chapter~3 establishes that business interruption is an operational/hazard risk, not a financial risk. This classification must be consistent across all later chapters.

\textbf{10-K Risk Factors as a pedagogical tool.} Chapter~3 introduces the SEC Form~10-K Risk Factors section as an analytical resource. Later chapters discussing specific risk types should cross-reference this framework and, where appropriate, include examples drawn from actual Risk Factors disclosures.

\begin{center}\rule{0.5\linewidth}{0.5pt}\end{center}

\begin{keytakeaways}
\item
  \textbf{One preamble, all chapters}: Every chapter uses \code{\textbackslash documentclass[11pt,openany]\{book\}} and \code{\textbackslash usepackage\{erm\_textbook\_style\}}. Do not modify the preamble template without updating this guide.
\item
  \textbf{Fixed opening sequence}: Learning objectives $\to$ chapter overview $\to$ corporate example $\to$ first section. This order is not negotiable.
\item
  \textbf{Eight box environments}: Each has a color and a purpose (Table~\ref{tab:boxes}). Using the wrong environment undermines the visual language readers learn to read.
\item
  \textbf{Minimum three sentences per paragraph}: One- and two-sentence paragraphs are prohibited in body text. Expand or merge.
\item
  \textbf{Named colors, never hex}: All colors come from the style package constants. The R theme uses the same palette, so figures match the page automatically.
  \item
  \textbf{Compile twice}: A single \code{pdflatex} pass will not resolve cross-references or TOC entries. Always run two passes before sharing a draft.
\end{keytakeaways}

\end{document}
